% =============================================================================
%  SPECTRA — Section IV: The SPECTRA Framework
%  Target: IEEE Transactions on Information Forensics & Security
%  Author: Sagar Korde
%
%  Required packages (add to preamble):
%    \usepackage{algorithm}
%    \usepackage{algpseudocode}
%    \usepackage{amsmath,amssymb,bm}
%    \usepackage{siunitx}
%    \usepackage{booktabs}
%    \usepackage{graphicx}
% =============================================================================

\section{The SPECTRA Framework}
\label{sec:framework}

\textsc{SPECTRA} is a seven-module pipeline that transforms a raw Bitcoin
transaction graph into per-address trust certificates, anomaly flags, and
multi-class transaction-type predictions.
Fig.~\ref{fig:architecture} gives a high-level view of the pipeline;
Algorithm~\ref{alg:spectra} provides the unified pseudocode.
Each module is designed to be independently replaceable while maintaining
well-defined input/output interfaces, enabling ablation and incremental
deployment in live blockchain-monitoring systems.
The modules are named by the letters of the \textsc{SPECTRA} acronym:
\textbf{S}pectral graph construction,
information-theoretic f\textbf{E}ature ranking,
\textbf{C}lustering via VAE and HDBSCAN,
Bayesian \textbf{P}robabilistic trust scoring,
Markov-chain \textbf{T}ransaction profiling,
GraphSAGE \textbf{R}ecognition, and
\textbf{A}nomaly and drift integration.
Table~\ref{tab:module_summary} summarizes the inputs, outputs, and key
hyperparameters of each module.

% ---------------------------------------------------------------------------
\begin{figure}[t]
  \centering
  % Replace the following rule with \includegraphics[width=\columnwidth]{figures/spectra_architecture.pdf}
  \rule{\columnwidth}{0.4pt}\\[2ex]
  \textit{[Architecture diagram: seven modules (S$\to$E$\to$C$\to$P$\to$T$\to$R$\to$A)
  shown as a left-to-right pipeline. Module~S receives the raw transaction
  graph and emits node embeddings; Module~E selects and ranks features;
  Module~C produces cluster assignments and latent codes; Modules~P and~T
  run in parallel to generate per-address trust scores and Markov anomaly
  scores respectively; Module~R fuses all representations for classification;
  Module~A issues anomaly flags and drift alerts. Data flows are depicted as
  directed arrows; feedback from Module~A to Module~C is shown as a dashed
  arc representing the Wasserstein drift monitor.]}\\[2ex]
  \rule{\columnwidth}{0.4pt}
  \caption{High-level architecture of \textsc{SPECTRA}.
    Seven modules process the Bitcoin transaction graph sequentially.
    Intermediate representations (spectral embeddings, VAE latent codes,
    cluster assignments, trust scores, and Markov anomaly scores) are
    concatenated to form the input feature vector for the GraphSAGE
    classifier (Module~R).
    The Wasserstein drift monitor in Module~A closes the feedback loop
    by triggering retraining alerts when the cluster-label distribution
    shifts beyond threshold $\delta=0.1$.}
  \label{fig:architecture}
\end{figure}
% ---------------------------------------------------------------------------

% ---------------------------------------------------------------------------
\subsection{Module S — Spectral Graph Construction}
\label{sec:module_s}

\paragraph{Graph construction.}
The input to \textsc{SPECTRA} is a directed weighted multi-graph
$\mathcal{G} = (\mathcal{V}, \mathcal{E}, \mathbf{W})$ derived from the
CryptoGraphNet dataset~\cite{cryptographnet}.
After parsing \num{300000} transactions we obtain
$|\mathcal{V}| = \num{2870000}$ address nodes and
$|\mathcal{E}| = \num{28600000}$ directed edges, where each edge weight
$w_{ij}$ equals the total BTC transferred from address $i$ to address $j$
across all co-occurring transactions.
The full graph is memory-intractable for spectral decomposition, so we prune
to the top-\num{30000} nodes ranked by weighted in-degree, retaining the
high-connectivity hub addresses that dominate the structural signal: exchange
hot wallets, mining pool payout addresses, and mixing-service intermediaries.
This degree-based pruning preserves the densest subgraph while reducing
computational cost by three orders of magnitude with negligible information
loss, since Bitcoin value flow is highly skewed~\cite{kondor2014rich}.

\paragraph{Normalized Laplacian.}
Let $\mathbf{A} \in \mathbb{R}^{N \times N}$ (with $N = \num{30000}$) be the
symmetrized adjacency matrix of the pruned subgraph and
$\mathbf{D} = \mathrm{diag}(d_1, \ldots, d_N)$ the corresponding degree
matrix.
We compute the symmetric normalized Laplacian
\begin{equation}
  \mathbf{L}_{\mathrm{sym}}
    = \mathbf{I} - \mathbf{D}^{-1/2}\mathbf{A}\mathbf{D}^{-1/2},
  \label{eq:lsym}
\end{equation}
whose eigenvalue spectrum is bounded in $[0, 2]$, facilitating stable
gradient flow in downstream modules.

\paragraph{Spectral embeddings via randomized SVD.}
We extract the $k = 10$ eigenvectors corresponding to the $k$ smallest
non-trivial eigenvalues of $\mathbf{L}_{\mathrm{sym}}$ as the spectral node
embedding $\mathbf{U} \in \mathbb{R}^{N \times k}$.
Classical ARPACK-based eigensolvers exhibit convergence instabilities on the
highly irregular degree distributions of real-world Bitcoin graphs (Zipfian
degree law with exponent $\approx 1.8$), manifesting as oscillating residuals
and occasional non-convergence beyond $k = 6$.
We therefore employ randomized SVD~\cite{halko2011finding} as implemented in
\texttt{sklearn.utils.extmath.randomized\_svd}, with
\texttt{n\_iter}$= 10$ power iterations and
\texttt{n\_oversamples}$= 20$ to control approximation error.
This yields a deterministic, convergent embedding at a fraction of the memory
footprint of a full dense eigendecomposition.

The spectral embedding $\mathbf{U}$ captures community structure: nodes
embedded close together in the spectral domain are highly likely to belong
to the same graph community~\cite{luxburg2007tutorial}, providing a
continuous cluster membership signal before any explicit clustering is
performed.

\paragraph{Behavioral feature augmentation.}
Six address-level behavioral features are appended to each spectral embedding:
total transaction count, mean transaction value, value standard deviation,
active lifespan in blocks, in-degree, and out-degree.
These features encode individual address patterns that are not captured by
community structure alone—a mixing service and an exchange may occupy similar
spectral positions yet differ markedly in transaction rhythm and value
variance.
The final node feature matrix is
$\mathbf{X}_{\mathrm{node}} \in \mathbb{R}^{N \times 16}$,
where $N = \num{30000}$ and the 16 dimensions comprise 10 spectral
components and 6 behavioral features.

% ---------------------------------------------------------------------------
\subsection{Module E — Information-Theoretic Feature Ranking}
\label{sec:module_e}

Module~E operates on the per-transaction feature matrix
$\mathbf{X}_{\mathrm{tx}} \in \mathbb{R}^{M \times F}$ (with
$M = \num{300000}$ transactions and $F$ raw features) and selects the most
informative feature subset for downstream learning, guarding against
redundancy and dimensionality inflation.

\paragraph{Shannon entropy screening.}
For each feature $X_j$ we compute the marginal Shannon entropy
\begin{equation}
  H(X_j) = -\sum_{x \in \mathcal{X}_j} p(x)\log_2 p(x),
  \label{eq:entropy}
\end{equation}
where $p(x)$ is estimated via histogram binning with Freedman--Diaconis
bandwidth.
Features with $H(X_j) < \epsilon_H$ (empirically $0.1$ bits) are discarded
as near-constant and carry no discriminative information.

\paragraph{Mutual information ranking.}
Mutual information between each surviving feature $X_j$ and the transaction
type label $Y$ is estimated as
\begin{equation}
  I(X_j;\, Y) = H(Y) - H(Y \mid X_j),
  \label{eq:mi}
\end{equation}
via the $k$-nearest-neighbor estimator~\cite{ross2014mutual} implemented in
\texttt{sklearn.feature\_selection.mutual\_info\_classif}.
Table~\ref{tab:mi_top20} lists the top-20 features by $I(X_j; Y)$.
Structural sizing features dominate: \texttt{weight} ($I = 1.320$),
\texttt{vsize} ($I = 1.204$), and \texttt{size} ($I = 1.193$) reflect
SegWit adoption patterns strongly correlated with transaction type.
Fee-related features (\texttt{fee\_rate\_sat\_per\_byte} $I = 0.494$,
\texttt{fee} $I = 0.471$) provide complementary signal tied to sender
urgency and mixing behavior.
We retain the top-20 features by $I(X_j; Y)$ as the transaction feature
vector $\mathbf{x}_{\mathrm{tx}} \in \mathbb{R}^{20}$.

\paragraph{NMF latent decomposition.}
To further compress the feature space and extract parts-based latent factors,
we apply Non-negative Matrix Factorization (NMF) with
$r = 10$ components, regularization coefficients
$\alpha_W = 0.1$ and $\ell_1$-ratio $= 0.5$
(elastic-net penalty on $\mathbf{W}$):
\begin{equation}
  \min_{\mathbf{W} \geq 0,\, \mathbf{H} \geq 0}
    \bigl\|\mathbf{X} - \mathbf{W}\mathbf{H}\bigr\|_F^2
    + \alpha_W \bigl[(1 - \ell_1)\tfrac{1}{2}\|\mathbf{W}\|_F^2
                    + \ell_1 \|\mathbf{W}\|_1\bigr].
  \label{eq:nmf}
\end{equation}
The Frobenius reconstruction error at convergence is $\num{1051.4}$,
confirming that $r = 10$ factors capture the dominant co-occurrence structure
in the top-20 feature block.
The NMF embedding $\mathbf{h} \in \mathbb{R}^{10}$ is concatenated with
$\mathbf{x}_{\mathrm{tx}}$ for input to Module~C.

% ---------------------------------------------------------------------------
\subsection{Module C — VAE and HDBSCAN Clustering}
\label{sec:module_c}

Module~C constructs a probabilistic latent space over transaction feature
vectors, enabling principled anomaly scoring and density-based clustering
that is robust to the class imbalance and irregular cluster shapes
characteristic of blockchain data.

\paragraph{Variational Autoencoder.}
The VAE encoder
$q_\phi(\mathbf{z} \mid \mathbf{x}) = \mathcal{N}(\boldsymbol{\mu}_\phi(\mathbf{x}),\,
\mathrm{diag}(\boldsymbol{\sigma}^2_\phi(\mathbf{x})))$
maps each transaction feature vector to a Gaussian distribution in latent
space $\mathcal{Z} = \mathbb{R}^{32}$.
The encoder uses fully connected layers with hidden dimensions $[128, 64]$
and dropout rate $0.2$; the decoder mirrors this architecture symmetrically.
The model is trained by maximizing the Evidence Lower BOund (ELBO):
\begin{equation}
  \mathcal{L}_{\mathrm{ELBO}}
    = \mathbb{E}_{q_\phi}\bigl[\log p_\theta(\mathbf{x} \mid \mathbf{z})\bigr]
      - \beta_{\mathrm{KL}}\,
        D_{\mathrm{KL}}\!\bigl(q_\phi(\mathbf{z} \mid \mathbf{x})
                               \;\|\; p(\mathbf{z})\bigr),
  \label{eq:elbo}
\end{equation}
with $\beta_{\mathrm{KL}} = 1.0$ (standard VAE) and isotropic Gaussian prior
$p(\mathbf{z}) = \mathcal{N}(\mathbf{0}, \mathbf{I})$.
Training proceeds for 50 epochs with batch size 512 and Adam
optimizer~\cite{kingma2014adam} (learning rate $0.001$).
Early stopping on the validation ELBO yields the best checkpoint at validation
loss $0.2216$.
The per-sample reconstruction error
$r(\mathbf{x}) = \|\mathbf{x} - \hat{\mathbf{x}}\|_2$
serves as the primary anomaly score in Module~A.

\paragraph{HDBSCAN density clustering.}
The VAE mean embedding $\boldsymbol{\mu}_\phi(\mathbf{x})$ is passed to
Hierarchical DBSCAN (HDBSCAN)~\cite{campello2013density} with
\texttt{min\_cluster\_size}$= 100$ and \texttt{min\_samples}$= 10$
under Euclidean metric.
HDBSCAN is preferred over $k$-means or Gaussian Mixture Models for three
reasons: (i) it discovers clusters of arbitrary shape in the latent manifold;
(ii) it explicitly models noise points rather than forcing them into the
nearest cluster; and (iii) its soft cluster membership probabilities are
available for downstream weighting.
The procedure yields \textbf{135 clusters} and
\num{63390} noise points (\SI{21.1}{\percent} of the corpus), with
normalized mutual information $\mathrm{NMI} = 0.112$ against the ground-truth
transaction-type labels—a meaningful signal given that clusters are learned
purely from graph topology and feature distributions, not label supervision.

\paragraph{Wasserstein temporal drift monitoring.}
To detect distributional shift in the transaction stream we partition
transactions into 30 consecutive monthly windows and compute the
Wasserstein-$1$ distance~\cite{villani2008optimal} between the cluster-label
distributions of adjacent windows:
\begin{equation}
  W_1(\mu_t, \mu_{t+1})
    = \inf_{\gamma \in \Gamma(\mu_t, \mu_{t+1})}
        \int \|c - c'\| \, d\gamma(c, c'),
  \label{eq:w1}
\end{equation}
where $\mu_t$ is the empirical cluster-assignment distribution in window $t$.
A drift alert is raised whenever $W_1(\mu_t, \mu_{t+1}) > \delta = 0.1$.
Across the 29 adjacent window pairs, all 29 pairs trigger the alert, with
$W_1$ scores ranging from $0.179$ to $50.745$, indicating substantial and
sustained distributional shift in Bitcoin transaction patterns over the
study period.
This persistent drift motivates SPECTRA's modular architecture: Module~C can
be periodically retrained in isolation without disrupting the other modules.

% ---------------------------------------------------------------------------
\subsection{Module P — Bayesian Trust Scoring}
\label{sec:module_p}

Module~P maintains a per-address trust score that is updated online as new
transactions are observed, providing a continuous and interpretable
credibility metric for each blockchain participant.

\paragraph{Beta-Bernoulli model.}
For each address $v$ we maintain a Beta posterior over the latent trust
parameter $\tau_v \in [0,1]$:
\begin{equation}
  p(\tau_v \mid \mathcal{D}_v)
    = \mathrm{Beta}\!\left(\alpha_v,\, \beta_v\right),
  \label{eq:beta_posterior}
\end{equation}
initialized with the weakly informative prior
$\mathrm{Beta}(2, 2)$ (mean $= 0.5$, mode at $0.5$, unimodal).
This prior encodes epistemic humility: absent any observed transactions, all
addresses are considered equally trustworthy.

\paragraph{Online posterior update.}
Upon observing transaction $\mathbf{x}$ at address $v$, the likelihood
of \emph{trustworthy behavior} is defined as
\begin{equation}
  \ell(\mathbf{x}) = \exp\!\bigl(-r(\mathbf{x})\bigr),
  \quad r(\mathbf{x}) = \|\mathbf{x} - \hat{\mathbf{x}}\|_2,
  \label{eq:likelihood}
\end{equation}
where $r(\mathbf{x})$ is the VAE reconstruction error from Module~C.
Low reconstruction error implies that the transaction conforms to the learned
distribution of typical activity for addresses in its cluster—high fidelity
behavior earns high trust.
The Beta parameters are updated via conjugate closed-form update:
\begin{equation}
  \alpha_v \leftarrow \alpha_v + \ell(\mathbf{x}), \quad
  \beta_v  \leftarrow \beta_v  + 1 - \ell(\mathbf{x}).
  \label{eq:beta_update}
\end{equation}
The posterior mean $\hat{\tau}_v = \alpha_v / (\alpha_v + \beta_v)$
is reported as the trust certificate $\tau_v$.

\paragraph{Population statistics.}
Across the \num{30000} hub addresses, the population trust distribution has
mean $\mu_\tau = 0.634$ and standard deviation $\sigma_\tau = 0.105$.
The right-skewed distribution—with a long low-trust tail—reflects the
Bitcoin ecosystem: a majority of addresses exhibit predictable, repetitive
behavioral patterns (exchange hot wallets reusing identical fee structures)
while a minority show highly variable, high-entropy activity consistent
with mixing or layering behavior.

% ---------------------------------------------------------------------------
\subsection{Module T — Markov Chain Behavioral Profiling}
\label{sec:module_t}

Module~T models the cluster-state trajectory of each address as a discrete
Markov chain, capturing behavioral patterns that unfold over sequences of
transactions rather than in individual events.

\paragraph{State space and transition estimation.}
The state space $\mathcal{S}$ consists of the $C = 135$ HDBSCAN cluster
labels plus one absorbing noise state, giving $|\mathcal{S}| = 136$.
From $\num{55678}$ multi-transaction address sequences we estimate the
maximum-likelihood transition matrix
$\mathbf{P} \in \mathbb{R}^{136 \times 136}$, where
\begin{equation}
  P_{ij} = \frac{n_{ij} + \varepsilon}{\sum_{j'} (n_{ij'} + \varepsilon)},
  \quad \varepsilon = 10^{-6},
  \label{eq:markov_transition}
\end{equation}
with $n_{ij}$ the observed count of transitions from state $i$ to state $j$
and $\varepsilon = 10^{-6}$ Laplace smoothing to avoid zero-probability
transitions for unobserved pairs.
The stationary distribution $\boldsymbol{\pi} = [\pi_1, \ldots, \pi_{136}]$
is computed via power iteration with tolerance $10^{-8}$.

\paragraph{Transition anomaly scoring.}
The anomaly score for a single transition $s_{t-1} \to s_t$ is the negative
log transition probability:
\begin{equation}
  \mathrm{score}(s_{t-1}, s_t) = -\log P_{s_{t-1}, s_t}.
  \label{eq:markov_score}
\end{equation}
High scores flag rare—and therefore anomalous—behavioral transitions, such as
an address that consistently transacts within one cluster (e.g., exchange
deposit cluster) suddenly appearing in the mixing-service cluster.
Across all observed transitions the population mean anomaly score is
$\bar{m} = 0.875$; Module~A flags an address if its per-transition score
exceeds $\bar{m} + 2\sigma_m$.

% ---------------------------------------------------------------------------
\subsection{Module R — GraphSAGE Transaction Recognition}
\label{sec:module_r}

Module~R performs multi-class transaction-type classification using a
two-layer GraphSAGE network~\cite{hamilton2017inductive} that aggregates
neighborhood information from the Bitcoin address graph.

\paragraph{Input feature concatenation.}
The input feature vector for address node $v$ is assembled as
\begin{equation}
  \mathbf{f}_v =
    \bigl[\,
      \mathbf{u}_v \;\|\;
      \boldsymbol{\mu}_v \;\|\;
      \tau_v \;\|\;
      m_v \;\|\;
      \mathbf{e}_v
    \,\bigr]
    \in \mathbb{R}^{d_f},
  \label{eq:feat_concat}
\end{equation}
where $\mathbf{u}_v \in \mathbb{R}^{16}$ is the spectral-behavioral embedding
from Module~S, $\boldsymbol{\mu}_v \in \mathbb{R}^{32}$ is the VAE latent
mean from Module~C, $\tau_v \in \mathbb{R}$ is the Bayesian trust score from
Module~P, $m_v \in \mathbb{R}$ is the Markov anomaly score from Module~T, and
$\mathbf{e}_v \in \mathbb{R}^{136}$ is the one-hot HDBSCAN cluster assignment
(including the noise label).
The total input dimension is $d_f = 16 + 32 + 1 + 1 + 136 = 186$.

\paragraph{GraphSAGE architecture.}
GraphSAGE~\cite{hamilton2017inductive} learns a neighborhood aggregation
function that generalizes inductively to unseen nodes:
\begin{equation}
  \mathbf{h}_v^{(l)}
    = \sigma\!\left(
        \mathbf{W}^{(l)}
        \cdot
        \mathrm{CONCAT}\!\left(
          \mathbf{h}_v^{(l-1)},\;
          \mathrm{MEAN}_{u \in \mathcal{N}(v)}\!\left\{
            \mathbf{h}_u^{(l-1)}
          \right\}
        \right)
      \right),
  \label{eq:graphsage}
\end{equation}
where $\sigma$ is the ReLU activation and $\mathcal{N}(v)$ is a fixed-size
sampled neighborhood.
We use two GraphSAGE layers with hidden dimension 128 and output embedding
dimension 64, followed by a linear classification head over six transaction
types.
Dropout with rate $0.3$ is applied after each layer.

\paragraph{Rationale for GraphSAGE over GAT.}
Graph Attention Networks (GAT)~\cite{velivckovic2018graph} learn attention
weights over neighbors but impose quadratic attention-head complexity per
node.
For the sparse, non-attributed Bitcoin address graph three considerations
favor GraphSAGE: (i) \emph{scalability} — neighborhood sampling in GraphSAGE
scales sub-linearly with graph size, critical for the \num{30000}-node pruned
subgraph and its highly skewed degree distribution; (ii) \emph{absence of
node attributes} — attention mechanisms derive most benefit from rich semantic
attributes (e.g., text or image features); on structural Bitcoin graphs the
gain is marginal while training cost increases significantly; and (iii)
\emph{inductive generalization} — GraphSAGE's mean aggregator generalizes to
unseen addresses appearing in the live transaction stream without retraining,
a critical operational requirement for real-time monitoring.

\paragraph{Training protocol.}
The model is trained for a maximum of 100 epochs using Adam with learning
rate $0.001$ and weight decay $10^{-4}$, with early stopping
(patience $= 15$ epochs) on the validation AUC-ROC.
Training terminates at epoch 18, yielding a total train time of
\SI{2.09}{\second} on the pruned subgraph—confirming that the modular
spectral preprocessing in Module~S has already extracted the dominant
structural signal, leaving a compact learning problem for the GNN.

% ---------------------------------------------------------------------------
\subsection{Module A — Anomaly and Drift Integration}
\label{sec:module_a}

Module~A is the decision layer of \textsc{SPECTRA}: it fuses the outputs of
Modules~C, P, and~T into per-address anomaly flags and drift alerts,
and emits a final trust certificate.

\paragraph{Unified anomaly flag.}
An address $v$ is flagged anomalous if \emph{any} of three complementary
conditions holds:
\begin{equation}
  \mathrm{anomaly}(v) = 1 \iff
    r(\mathbf{x}_v) > \tau_{\mathrm{anom}}
    \;\lor\;
    m_v > \bar{m} + 2\sigma_m
    \;\lor\;
    \mathcal{W}(t) > \delta,
  \label{eq:anomaly_flag}
\end{equation}
where $r(\mathbf{x}_v)$ is the VAE reconstruction error,
$\tau_{\mathrm{anom}} = 0.1694$ is the $95^{\mathrm{th}}$ percentile of the
population reconstruction error (learned from the training set),
$m_v$ is the Markov anomaly score from \eqref{eq:markov_score},
$\bar{m}$ and $\sigma_m$ are its population mean and standard deviation, and
$\mathcal{W}(t) = W_1(\mu_{t-1}, \mu_t)$ is the Wasserstein distance for
time window $t$ with threshold $\delta = 0.1$.
The disjunctive combination ensures that an address is flagged by any
evidence type; this union-of-detectors design minimizes false negatives at the
cost of a controlled increase in false positives, which is the preferred
operating point for financial forensics~\cite{bahnsen2016fraud}.

\paragraph{Trust certificate.}
The final trust certificate for address $v$ is the posterior mean trust score
$\hat{\tau}_v = \alpha_v / (\alpha_v + \beta_v) \in [0,1]$ from Module~P,
optionally attenuated by the anomaly flag:
$\tilde{\tau}_v = \hat{\tau}_v \cdot (1 - \mathrm{anomaly}(v))$.
This scalar certificate is the primary output delivered to downstream
compliance or forensic systems.

\paragraph{Drift alert.}
In addition to per-address flags, Module~A emits a binary per-window drift
alert $\mathcal{A}_t = \mathbf{1}[\mathcal{W}(t) > \delta]$.
Sustained drift alerts over consecutive windows (all 29 of the 29 observable
pairs in the study period) indicate concept drift requiring retraining of
Module~C's VAE and HDBSCAN components.
Because the Markov transition matrix (Module~T) and GraphSAGE weights
(Module~R) are conditioned on the cluster state space defined by Module~C,
a cluster re-labeling event requires propagating updated cluster assignments
through Modules~P, T, and~R; the modular architecture ensures this can be
done without restarting the pipeline.

% ---------------------------------------------------------------------------
\subsection{Unified Algorithm}
\label{sec:algorithm}

Algorithm~\ref{alg:spectra} presents the complete \textsc{SPECTRA} pipeline
in pseudocode, showing the sequential data flow across all seven modules.
Notation follows~\eqref{eq:lsym}--\eqref{eq:anomaly_flag}.

% ---------------------------------------------------------------------------
\begin{algorithm}[t]
\caption{\textsc{SPECTRA} — Spectral-Probabilistic Ensemble Clustering for
         Transaction Recognition and Authentication}
\label{alg:spectra}
\begin{algorithmic}[1]
\Require Raw transaction dataset $\mathcal{D}$ (\num{300000} transactions);
         hyperparameters $k, r, \beta_{\mathrm{KL}}, \delta, \varepsilon$
\Ensure  Per-address trust certificate $\tilde{\tau}_v$,
         anomaly flag $\mathrm{anomaly}(v)$,
         transaction-type prediction $\hat{y}_v$,
         drift alert $\mathcal{A}_t$

\Statex
\Statex \textbf{Module S — Spectral Graph Construction}
\State Build directed weighted address graph
       $\mathcal{G}=(\mathcal{V},\mathcal{E},\mathbf{W})$
       from $\mathcal{D}$\quad
       ($|\mathcal{V}|=\num{2870000}$, $|\mathcal{E}|=\num{28600000}$)
\State Prune $\mathcal{G}$ to top-\num{30000} nodes by weighted in-degree
       $\Rightarrow \mathcal{G}'$
\State Compute symmetric normalized Laplacian
       $\mathbf{L}_{\mathrm{sym}} = \mathbf{I} -
        \mathbf{D}^{-1/2}\mathbf{A}\mathbf{D}^{-1/2}$
       on $\mathcal{G}'$
\State Extract $k=10$ spectral embeddings
       $\mathbf{U} \in \mathbb{R}^{N \times 10}$
       via randomized SVD (\texttt{n\_iter}$=10$, \texttt{n\_oversamples}$=20$)
\State Compute 6 behavioral features per node; concatenate with $\mathbf{U}$
       $\Rightarrow \mathbf{X}_{\mathrm{node}} \in \mathbb{R}^{N \times 16}$

\Statex
\Statex \textbf{Module E — Information-Theoretic Feature Ranking}
\State Compute Shannon entropy $H(X_j)$ for each raw feature $X_j$;
       discard low-entropy features
\State Estimate $I(X_j; Y)$ via $k$-NN mutual information;
       retain top-20 features
       $\Rightarrow \mathbf{x}_{\mathrm{tx}} \in \mathbb{R}^{M \times 20}$
\State Fit NMF ($r=10$, $\alpha_W=0.1$, $\ell_1$-ratio$=0.5$)
       $\Rightarrow \mathbf{H} \in \mathbb{R}^{M \times 10}$;
       concatenate $[\mathbf{x}_{\mathrm{tx}} \| \mathbf{H}]$

\Statex
\Statex \textbf{Module C — VAE + HDBSCAN Clustering}
\State Train VAE (hidden$=[128,64]$, latent$=32$, $\beta_{\mathrm{KL}}=1.0$,
       dropout$=0.2$, 50 epochs, lr$=0.001$) on $\mathbf{x}_{\mathrm{tx}}$
\State Encode all transactions:
       $(\boldsymbol{\mu}_v, \boldsymbol{\sigma}^2_v)
        \leftarrow q_\phi(\mathbf{z} \mid \mathbf{x}_v)$;
       record reconstruction error $r(\mathbf{x}_v)$
\State Cluster $\{\boldsymbol{\mu}_v\}$ with HDBSCAN
       (\texttt{min\_cluster\_size}$=100$, \texttt{min\_samples}$=10$)
       $\Rightarrow$ 135 clusters + noise; labels $\{c_v\}$
\For{each adjacent monthly window pair $(t-1, t)$, $t=1,\ldots,29$}
    \State Compute $\mathcal{W}(t) = W_1(\mu_{t-1}, \mu_t)$;
           set $\mathcal{A}_t = \mathbf{1}[\mathcal{W}(t) > \delta]$
\EndFor

\Statex
\Statex \textbf{Module P — Bayesian Trust Scoring}
\State Initialize $(\alpha_v, \beta_v) \leftarrow (2, 2)$ for all $v$
\For{each transaction $\mathbf{x}$ at address $v$ (in temporal order)}
    \State $\ell \leftarrow \exp(-r(\mathbf{x}))$
    \State $\alpha_v \leftarrow \alpha_v + \ell$;\quad
           $\beta_v \leftarrow \beta_v + (1-\ell)$
\EndFor
\State $\hat{\tau}_v \leftarrow \alpha_v / (\alpha_v + \beta_v)$
       for all $v$

\Statex
\Statex \textbf{Module T — Markov Chain Behavioral Profiling}
\State Build transition count matrix $\mathbf{N} \in \mathbb{R}^{136 \times 136}$
       from address cluster-state sequences
\State Estimate $\mathbf{P}$:
       $P_{ij} \leftarrow (n_{ij} + \varepsilon) /
                           (\sum_{j'} n_{ij'} + 136\varepsilon)$,\;
       $\varepsilon = 10^{-6}$
\State Compute stationary distribution $\boldsymbol{\pi}$ via power iteration
       (tol $=10^{-8}$)
\State $m_v \leftarrow -\log P_{c_{v,t-1},\, c_{v,t}}$
       for each transition of address $v$

\Statex
\Statex \textbf{Module R — GraphSAGE Transaction Recognition}
\State Assemble node feature matrix:
       $\mathbf{f}_v =
        [\mathbf{X}_{\mathrm{node},v} \;\|\;
         \boldsymbol{\mu}_v \;\|\;
         \hat{\tau}_v \;\|\;
         m_v \;\|\;
         \mathbf{e}_{c_v}]
        \in \mathbb{R}^{186}$
\State Train two-layer GraphSAGE (hidden$=128$, embed$=64$, dropout$=0.3$,
       lr$=0.001$, wd$=10^{-4}$, patience$=15$) on $\mathcal{G}'$
       with node features $\mathbf{F}$ and transaction-type labels
\State Infer $\hat{y}_v \leftarrow \mathrm{argmax}\;\mathrm{softmax}
       (\mathbf{h}_v^{(2)}\mathbf{W}_{\mathrm{cls}})$ for all $v$

\Statex
\Statex \textbf{Module A — Anomaly and Drift Integration}
\State Compute population statistics:
       $\tau_{\mathrm{anom}} \leftarrow \mathrm{perc}_{95}(\{r(\mathbf{x}_v)\})$;\quad
       $\bar{m}, \sigma_m \leftarrow \mathrm{mean}, \mathrm{std}(\{m_v\})$
\For{each address $v$}
    \State $\mathrm{anomaly}(v) \leftarrow
            \mathbf{1}\!\left[
              r(\mathbf{x}_v) > \tau_{\mathrm{anom}}
              \;\lor\;
              m_v > \bar{m} + 2\sigma_m
              \;\lor\;
              \mathcal{A}_{t(v)} = 1
            \right]$
    \State $\tilde{\tau}_v \leftarrow \hat{\tau}_v \cdot
            (1 - \mathrm{anomaly}(v))$
\EndFor
\State \Return $\{\hat{y}_v\}$, $\{\tilde{\tau}_v\}$,
               $\{\mathrm{anomaly}(v)\}$, $\{\mathcal{A}_t\}$
\end{algorithmic}
\end{algorithm}
% ---------------------------------------------------------------------------

\subsection{Complexity Analysis}
\label{sec:complexity}

Table~\ref{tab:complexity} summarizes the computational complexity of each
module.
The dominant cost is Module~S: randomized SVD on $\mathbf{L}_{\mathrm{sym}}$
runs in $\mathcal{O}(N k \log k)$ time~\cite{halko2011finding}, far cheaper
than the $\mathcal{O}(N^2)$ dense eigensolver.
Module~C's VAE training is $\mathcal{O}(M \cdot B^{-1} \cdot E_{\max})$
where $B = 512$ is the batch size and $E_{\max} = 50$ the epoch budget;
HDBSCAN is $\mathcal{O}(M \log M)$ for the mutual reachability graph
construction phase.
Module~T's transition matrix estimation is $\mathcal{O}(|\mathcal{S}|^2)$
for matrix storage and $\mathcal{O}(|\mathcal{S}|^2 / (1-\lambda_2))$ for
power iteration convergence, where $\lambda_2$ is the second eigenvalue of
$\mathbf{P}$.
Module~R's GraphSAGE forward pass is $\mathcal{O}(N \cdot d_f \cdot L \cdot S)$
where $L = 2$ is the number of layers and $S = 25$ the neighborhood sample
size; early stopping at epoch 18 further limits wall-clock training time to
\SI{2.09}{\second}.
End-to-end inference for a new address (given frozen modules) is dominated by
the VAE encoder ($\mathcal{O}(d_f \cdot 128)$) and GraphSAGE forward pass
($\mathcal{O}(d_f \cdot S \cdot 128)$), both operating in sub-millisecond
wall-clock time on a single CPU core.

\begin{table}[t]
  \caption{Per-module complexity and key hyperparameters of \textsc{SPECTRA}.
           $N=\num{30000}$ (pruned nodes), $M=\num{300000}$ (transactions),
           $k=10$ (spectral dims), $C=136$ (cluster states),
           $d_f=186$ (GNN input dim).}
  \label{tab:module_summary}
  \centering
  \renewcommand{\arraystretch}{1.15}
  \begin{tabular}{@{}lllp{3.6cm}@{}}
    \toprule
    \textbf{Module} & \textbf{Output} & \textbf{Complexity} & \textbf{Key hyperparameters} \\
    \midrule
    S — Spectral  & $\mathbf{X}_{\mathrm{node}} \!\in\! \mathbb{R}^{N\times 16}$
                  & $\mathcal{O}(Nk\log k)$
                  & $k{=}10$, \texttt{n\_iter}$=10$ \\
    E — Entropy   & $\mathbf{x}_{\mathrm{tx}} \!\in\! \mathbb{R}^{M\times 20}$
                  & $\mathcal{O}(MF)$
                  & top-20 MI, NMF $r{=}10$ \\
    C — Cluster   & $\{c_v\}$, $\{\boldsymbol{\mu}_v\}$, $\{\mathcal{A}_t\}$
                  & $\mathcal{O}(M\log M)$
                  & latent$=32$, MCS$=100$ \\
    P — Probabilistic & $\hat{\tau}_v \!\in\! [0,1]$
                  & $\mathcal{O}(M)$
                  & $\mathrm{Beta}(2,2)$ prior \\
    T — Markov    & $m_v \!\in\! \mathbb{R}_{+}$
                  & $\mathcal{O}(C^2)$
                  & $|\mathcal{S}|{=}136$, $\varepsilon{=}10^{-6}$ \\
    R — Recognition & $\hat{y}_v \!\in\! \{1,\ldots,6\}$
                  & $\mathcal{O}(NLSd_f)$
                  & hidden$=128$, dropout$=0.3$ \\
    A — Anomaly   & $\mathrm{anomaly}(v)$, $\tilde{\tau}_v$
                  & $\mathcal{O}(N)$
                  & $\tau_{\mathrm{anom}}{=}0.1694$, $\delta{=}0.1$ \\
    \bottomrule
  \end{tabular}
\end{table}

\begin{table}[t]
  \caption{Top-20 transaction features ranked by mutual information
           $I(X_j; Y)$ with transaction-type label $Y$.
           All values rounded to three decimal places.}
  \label{tab:mi_top20}
  \centering
  \renewcommand{\arraystretch}{1.12}
  \begin{tabular}{@{}clc@{}}
    \toprule
    \textbf{Rank} & \textbf{Feature} & $I(X_j;\,Y)$ \\
    \midrule
    1  & \texttt{weight}                   & 1.320 \\
    2  & \texttt{vsize}                    & 1.204 \\
    3  & \texttt{size}                     & 1.193 \\
    4  & \texttt{output\_count}            & 1.171 \\
    5  & \texttt{input\_count}             & 0.726 \\
    6  & \texttt{fee\_rate\_sat\_per\_byte}& 0.494 \\
    7  & \texttt{fee}                      & 0.471 \\
    8  & \texttt{value\_difference}        & 0.470 \\
    9  & \texttt{fee\_rate\_sat\_per\_vbyte}& 0.410 \\
    10 & \texttt{total\_output\_value}     & 0.380 \\
    \multicolumn{3}{c}{\textit{(ranks 11--20 reported in supplementary material)}} \\
    \bottomrule
  \end{tabular}
\end{table}
